\documentclass{article}
\usepackage{algorithm, algorithmic, amsfonts, amsmath}

\begin{document}
\begin{algorithm}[!h]
\renewcommand{\algorithmicrequire}{\textbf{Input:}}
\renewcommand{\algorithmicensure}{\textbf{Output:}}
	\caption{PURE: positive-unlabeled recommendation with a generative adversarial network}
	\begin{algorithmic}[1]
	    \REQUIRE
	    Generators $G_{u},G_{i}$; discriminator $D$; user--item interaction matrix $\mathcal{R}$;
	    user set $\mathcal{U}$; item set $\mathcal{I}$; positive class prior $\pi_{p}$
	    \ENSURE
	    The trained $G_{u}$, $G_{i}$ and $D$
	    \STATE Assign $G_{u},G_{i}$ random weights; assign $D$ random or pre-trained weights;
	    $n_{p}=|\Omega|$,\quad $n_{u}=\mathrm{ceil}\!\left(\dfrac{n_{p}}{(1-2\pi_{p})^{2}}\right)$ % the unlabeled sample size follows from the prior
	    \REPEAT
	    \FOR{discriminator-steps}
	    \STATE Sample first $n_{p}$ tuples $(u,i)\in\Omega$ with label $1$
	    \STATE Sample another $n_{p}$ tuples $(u,i)\in\Omega$ with label $0$ % the same observed positives enter the objective under both labels, which is the PU risk correction
	    \STATE Sample $n_{u}$ tuples $(u,i)\in\Omega^{-}$ with label $0$
	    \STATE Generate $n_{u}$ tuples $(u,i')$ and $(u',i)$ with label $0$ using Eq.~(13)
	    \STATE Update $D$ by ascending its gradients in the objectives of Eq.~(10) and Eq.~(11):
	    \STATE \quad $V(D)_{1}=\sum_{(u,i)\in\Omega}\pi_{p}\log D(u,i)-\pi_{p}\log\!\left(1-D(u,i)\right)$
	    \STATE \quad $V(D)_{2}=\sum_{(u,i)\in\Omega^{-}}\log\!\left(1-D(u,i)\right)
	    +\left[\log\!\left(1-D(u,i')\right)+\log\!\left(1-D(u',i)\right)\right]$
	    \ENDFOR
	    \FOR{generator-steps}
	    \STATE Generate $n_{u}$ random noise $z_{u},z_{i}\sim\mathcal{N}(\mathbf{0},\delta I)$ (Eq.~12)
	    \STATE Sample $n_{u}$ tuples $(u,i)\in\Omega^{-}$ with label $1$
	    \STATE Replace $(u,i)$ with $(u,i')$ and $(u',i)$ using $G_{u}(z_{u})$, $G_{i}(z_{i})$ (Eq.~13,
	    a two-layer ReLU MLP) % the generated users and items are virtual and need not exist in I
	    \STATE Update $G_{u},G_{i}$ by descending their gradients in the objective of Eq.~(11)
	    \ENDFOR
	    \UNTIL{convergence}
		\RETURN $G_{u}$, $G_{i}$, $D$
	\end{algorithmic}
\end{algorithm}
% Zhou, Xu, Wu, Wang, He, KDD 2021, "PURE: Positive-Unlabeled Recommendation with
% Generative Adversarial Network". Transcribed line by line from Algorithm 1 and equations (10) to (14).
% Setting: implicit-feedback recommendation is naturally PU. An item a user clicked is a positive, and one not clicked
% is merely unobserved, not disliked. The PU correction sits at lines 4 and 5 and equation (10): the observed
% positives Omega enter both with label 1 and, weighted by pi_p, with label 0. The second term is the empirical risk
% of the positives against the negative label, and it cancels the bias from treating every unobserved item as negative.
% The generator differs in role from PGAN: it produces virtual users u' and items i', which need not exist in I, as
% additional unlabeled samples for the discriminator objective, not as substitutes for negatives. The unbiasedness of
% the PU risk is therefore guaranteed by the prior weighting of equation (10) and is decoupled from generation quality.
% The unlabeled sample size n_u = ceil(n_p / (1-2 pi_p)^2) follows directly from the class prior. The original fixes
% the ratio of generated to sampled unlabeled instances at 1 and notes that tuning it may improve results further.
\end{document}
